\section{Introduction}
\label{sec:introduction}

All-solid-state lithium batteries require an electrolyte that combines rapid lithium-ion transport with mechanical integrity, chemical tolerance and a manufacturable form factor. Garnet-type Li$_7$La$_3$Zr$_2$O$_{12}$ (LLZO) is attractive because its oxide framework supports a wide range of aliovalent substitutions and can be processed as dense ceramics, thin sheets or composite electrodes \cite{wang2020garnet,puente2021garnet,kim2020research}. Yet the name LLZO covers samples made with different lithium inventories, precursor chemistries, atmospheres, covers, crucibles, thermal ramps and impedance geometries. A conductivity number therefore cannot be interpreted without the material state and the route that produced it.

This review follows a process--structure--transport chain. It asks how composition and phase stability set the available lithium sites; how precursor chemistry and forming determine the green body; how thermal history and lithium activity set density and grain-boundary chemistry; and how those states are read by impedance, spectroscopy and cell tests. The resulting comparison is deliberately conditional: a route is considered transferable only when the composition, thermal history, relative density, microstructure and measurement definition are reported together. The framework extends the processing emphasis of recent LLZO reviews \cite{cao2019modeling,wang2024accelerating,zhao2024review,sazvar2025review} by making comparability itself an object of analysis.

\subsection{Scope and evidence criteria}

The target phase is garnet-type LLZO and its Al-, Ta-, Ga-, Nb-, Y-, Sr- and related substituted variants. We include powder synthesis, calcination, forming, sintering, rapid or assisted densification, membranes, grain boundaries, environmental exposure, Li contacts and cathode integration. We use related garnets as mechanistic comparators, but do not transfer their conductivity values to LLZO. We also separate direct measurements from calculations and from nominal-composition assignments. Reviews provide vocabulary and route families, while primary studies anchor claims about processing or electrochemical response \cite{wang2020garnet,puente2021garnet,raju2021crystal,chen2022processing,wang2026research}.

The scope is also informed by a deliberately bounded oxide-ceramic analogy. In BaZn$_2$Si$_2$O$_7$, diffraction and thermal-expansion work resolves a temperature-triggered structural transition, while Ba/Sr substitution stabilizes a related high-temperature framework \cite{lin1999phase,thieme2015ba1}; combustion and crystal-growth studies show that precursor route and dopant choice change phase selection and luminescence \cite{yao2011synthesis,yao2011crystal}. Partial BaO--SrO--ZnO--SiO$_2$ and BaO--ZnO--SiO$_2$--B$_2$O$_3$ studies further illustrate how additive content and eutectic liquid formation can redirect phase assemblages \cite{thieme2018effect,chen2010phase}. Thermodynamic treatments of Y$_2$O$_3$--SiO$_2$, ZnO--SiO$_2$ and related oxide subsystems provide a methodological precedent for combining measured equilibria with CALPHAD extrapolation \cite{makrovets2023y2o3,shevchenko2019thermodynamic,konar2020a}. No exact Ba--Y--Zn--Si--O quaternary phase diagram or thermodynamic assessment was identified in the bounded corpus; these partial-system results therefore cannot establish Ba--Y--Zn--Si--O phase fields. The silicate papers are not LLZO measurements and do not establish LLZO phase fields or conductivity; they are used only to motivate an explicit workflow of phase-diagram measurement, composition spacing and thermodynamic modeling.

The minimum record for a meaningful comparison contains (i) nominal and measured composition, including dopant and lithium excess; (ii) precursor identity, mixing and calcination; (iii) pellet or tape geometry, cover powder, atmosphere and crucible; (iv) sintering schedule and relative density; (v) phase fractions, grain size, porosity and boundary chemistry; and (vi) conductivity type, temperature, electrode configuration, fitting model and applied pressure. Missing items do not invalidate a study, but they limit the strength of a cross-paper comparison. This rule is consistent with the measurement cautions emphasized for garnet electrolytes and critical-current tests \cite{zhao2024review,wang2024review,klimpel2023standardizing}.

\subsection{Near-neighbor garnet systems}
\label{sec:near-neighbor}

Before focusing on LLZO, it is useful to state what neighboring garnets have already taught us. Li$_5$La$_3$Ta$_2$O$_{12}$ and other lithium-stuffed garnets show that cation substitution changes lithium concentration, site occupancy and the tetragonal--cubic transition, while pressure-assisted consolidation can alter density independently of composition \cite{zhao2018enhancement,cao2019modeling,puente2021garnet}. These results support a transferable principle: phase symmetry and carrier concentration are coupled, but neither predicts the measured total resistance without a microstructural descriptor.

The same family also demonstrates why processing details matter. Thin-sheet, water-based and self-standing membrane studies report dense garnet bodies at thermal budgets different from conventional LLZO firing, and their transport depends on thickness, residual porosity and electrode geometry \cite{gao2019preparation,ye2020water,zhang2023ultrafast}. Atomistic and microstructural studies further indicate that low-energy pathways can coexist with mechanically vulnerable interfaces \cite{heo2021microstructural,yu2018grain}. We therefore use neighboring garnets as evidence for mechanisms and reporting practices, not as a ranking of LLZO recipes.

Foundational surveys and application-oriented studies converge on the same reporting problem: garnet composition, processing, interfaces and cell architecture are inseparable when a result is transferred across laboratories \cite{cheng2024li,kodgire2024review,han2023recent,liu2018challenges,jia2020comprehensive,li2020status,sun2022prospects,guo2022solid}. Recent processing guidance adds scale-up, environmental and manufacturing constraints to that comparison \cite{balaish2025emerging,touidjine2026tapex,jagan2026recent}. These sources frame the review, whereas the mechanistic claims below are tied to primary LLZO or explicitly marked comparator studies.

\subsection{Roadmap}

Figure~\ref{fig:roadmap} is a reading map. Section~\ref{sec:background} defines phase, defect and transport terms; Sections~\ref{sec:composition}--\ref{sec:densification} follow the synthesis chain; Section~\ref{sec:transport} reconciles conductivity measurements; Section~\ref{sec:interfaces} treats Li and cathode contacts; and Sections~\ref{sec:discussion}--\ref{sec:open} identify discriminating experiments. The process--structure--conductivity map is used in Section~\ref{sec:composition}, while the coupled-variable feedback map is interpreted in Section~\ref{sec:discussion}.

\begin{figure}[t]
  \centering
  \includegraphics[width=0.93\linewidth]{../figures/svg/fig2_reading_roadmap_preview.png}
  \caption{Reading roadmap for the review. The sequence moves from composition and phase state to powder, densification, microstructure, transport measurement and cell integration.}
  \label{fig:roadmap}
\end{figure}
